\begin{helper}
Let \(N\) be a natural number, and let \(c, b, x \in \mathbb R^{\operatorname{Fin}(N)}\) have nonnegative integer entries and equal total sums. Suppose \(b\) is nonincreasing, \(x\) is a nonincreasing permutation of \(c\), and for every \(0 \le r \le N\), the prefix sum of \(x\) is at least the prefix sum of \(b\). If \(h \colon \mathbb R^{\operatorname{Fin}(N)} \to \mathbb R\) is convex and symmetric under coordinate permutations, then \(h(b) \le h(x)\).
\end{helper}
\begin{proof}
Use well-founded induction on the integer-valued distance
\[
D(u)=\sum_{k=1}^N |u_k-b_k|
\]
among vectors \(u\) that are nonincreasing, have nonnegative integer entries,
have the same total sum as \(b\), and dominate \(b\) on all prefixes.  This
is an integer measure because every entry of \(u\) and \(b\) is a
nonnegative integer, so every summand is an absolute value of an integer
difference.  If \(u=b\), then \(h(b)\le h(u)\) is equality.  It remains to
show that every \(u\ne b\) in this induction domain admits a new vector in
the same domain with strictly smaller \(D\), and with \(h\) no larger than
at \(u\).  We prove this for the given \(x\).

Assume \(x\ne b\).  By \(\noderef{MajorizationDiscrepancyTransfer}\), there
are indices \(i<j\) and a positive integer \(\delta\) such that
\[
b_i<x_i,\qquad x_j<b_j,\qquad
\delta\le x_i-b_i,\qquad \delta\le b_j-x_j,
\]
and for every intermediate prefix \(i<r\le j\), the prefix slack of \(x\)
over \(b\) is at least \(\delta\).  Define \(y\) by moving \(\delta\) units
from coordinate \(i\) to coordinate \(j\):
\[
y_i=x_i-\delta,\qquad y_j=x_j+\delta,\qquad y_k=x_k\quad(k\ne i,j).
\]
The vector \(y\) has nonnegative integer entries.  Indeed, the entries away
from \(i,j\) are unchanged; \(y_j\) is the sum of two integers; and
\(\delta\le x_i-b_i\) with \(b_i\ge0\) gives \(x_i-\delta\ge b_i\ge0\).
By \(\noderef{MajorizationTransferPrefixDominance}\), \(y\) has the same
total sum as \(b\) and still dominates \(b\) on every prefix.

The same transfer strictly lowers the distance before resorting.  At
coordinate \(i\), the inequalities \(b_i<x_i\) and
\(\delta\le x_i-b_i\) give
\[
|y_i-b_i|=|x_i-\delta-b_i|=x_i-b_i-\delta
       =|x_i-b_i|-\delta.
\]
At coordinate \(j\), the inequalities \(x_j<b_j\) and
\(\delta\le b_j-x_j\) give
\[
|y_j-b_j|=|x_j+\delta-b_j|=b_j-x_j-\delta
       =|x_j-b_j|-\delta.
\]
Every other coordinate contribution to \(D\) is unchanged, so
\[
D(y)=D(x)-2\delta<D(x).
\]
By \(\noderef{MajorizationTransferConvexCombination}\), \(y\) is a convex
combination
\[
y=(1-\theta)x+\theta\,\tau(x)
\]
for some \(0\le\theta\le1\), where \(\tau\) is the transposition swapping
\(i\) and \(j\).  Convexity of \(h\) gives
\[
h(y)\le (1-\theta)h(x)+\theta h(\tau(x)).
\]
Symmetry gives \(h(\tau(x))=h(x)\), hence \(h(y)\le h(x)\).

Choose a nonincreasing rearrangement \(y^\downarrow\) of \(y\).  We next
prove the prefix-maximum property needed to keep prefix dominance after
resorting.  Let \(F_r\) be the first \(r\) positions.  For any set \(T\) of
\(r\) positions in a nonincreasing vector, the sum over \(F_r\) is at least
the sum over \(T\).  If \(T\ne F_r\), choose
\(t\in T\setminus F_r\).  Since \(|T|=|F_r|=r\), there is also some
\(s\in F_r\setminus T\).  Then \(s<t\), so nonincreasing order gives
\(y^\downarrow_s\ge y^\downarrow_t\).  Replacing \(t\) by \(s\) therefore
does not decrease the selected sum.  Repeating this finite exchange removes
all elements outside \(F_r\), and proves
\[
\sum_{k<r} y^\downarrow_k
\ge
\sum_{t\in T} y^\downarrow_t
\]
for every \(r\)-element set \(T\).

Apply this to the \(r\) positions of \(y^\downarrow\) that correspond, under
a permutation witnessing that \(y^\downarrow\) is a rearrangement of \(y\),
to the original first \(r\) coordinates of \(y\).  The selected sum is
exactly \(\sum_{k<r}y_k\).  Thus
\[
\sum_{k<r} y^\downarrow_k\ge \sum_{k<r} y_k.
\]
Since \(y\) already dominates \(b\) on prefixes, it follows that
\[
\sum_{k<r} b_k\le \sum_{k<r}y_k\le \sum_{k<r}y^\downarrow_k
\]
for every \(0\le r\le N\).  Hence \(y^\downarrow\) dominates \(b\) on all
prefixes.

We now check that \(y^\downarrow\) is again in the induction domain.  It is
nonincreasing by construction.  It has nonnegative integer entries because
it is a permutation of the nonnegative-integer vector \(y\).  It has the
same total sum as \(y\), hence the same total sum as \(b\).  The preceding
paragraph proves the required prefix dominance.  Finally,
\(\noderef{MajorizationResortingDistance}\) gives
\[
D(y^\downarrow)\le D(y)<D(x),
\]
because \(b\) and \(y^\downarrow\) are nonincreasing and
\(y^\downarrow\) is a permutation of \(y\).

Symmetry also gives \(h(y^\downarrow)=h(y)\), so
\[
h(y^\downarrow)\le h(x).
\]
By the induction hypothesis applied to \(y^\downarrow\), whose distance from
\(b\) is strictly smaller and which satisfies all induction-domain
hypotheses, we have
\[
h(b)\le h(y^\downarrow).
\]
Combining the last two inequalities yields \(h(b)\le h(x)\).  This completes
the induction and proves the claim.
\end{proof}
